\documentclass[../DoAn.tex]{subfiles}
\begin{document}

% ═══════════════════════════════════════════════════════════════════════════════
\section{Thiết kế Kiến trúc Tổng thể}
\label{sec:architecture}
% ═══════════════════════════════════════════════════════════════════════════════

Hệ thống được xây dựng theo mô hình ứng dụng phi tập trung (Decentralized Application -- DApp) với kiến trúc lai ba tầng giữa thành phần on-chain, off-chain và client, nhằm cân bằng giữa tính phi tập trung, hiệu năng xử lý và trải nghiệm người dùng~\cite{erc4337,groth16}.

\begin{figure}[H]
    \centering
    \includegraphics[width=1\textwidth]{SOICT_DATN_NGUYEN_VAN_AN_20224910 (1)/Hinhve/overall_architecture.png}
    \caption{Kiến trúc tổng thể ba tầng của hệ thống Stealth Wallet}
    \label{fig:ArchitectureDiagram}
\end{figure}

\subsection{Tầng Smart Contract (On-chain Layer)}
\label{subsec:onchain_layer}

Tầng này đảm nhiệm việc lưu trữ trạng thái không thể thay đổi và xác minh các bằng chứng mật mã trên blockchain Ethereum Sepolia. Các hợp đồng thông minh được triển khai bao gồm:

\begin{itemize}
    \item \textbf{\texttt{StealthTreeManager}:} Lưu trữ \texttt{Merkle Root} hiện hành của cây Merkle Poseidon độ sâu 20. Hợp đồng tích hợp trực tiếp với \texttt{SMTUpdateVerifier} để xác minh bằng chứng Groth16 trước mỗi lần cập nhật gốc cây. Đây là điểm đồng thuận duy nhất giữa trạng thái off-chain (server) và on-chain.

    \item \textbf{\texttt{StealthAccountFactory}:} Factory triển khai các \texttt{StealthAccount} thông qua cơ chế \texttt{CREATE2}. Địa chỉ stealth được xác định trước dựa trên \texttt{indexCommitment} đóng vai trò là \texttt{salt}, cho phép tài sản gửi đến tài khoản trước khi hợp đồng được triển khai thực sự~\cite{eip1014}.

    \item \textbf{\texttt{StealthAccount}:} Smart Account tương thích ERC-4337, tích hợp:
    \begin{itemize}
        \item Hàm \texttt{validateUserOp()}: xác thực \texttt{UserOperation} thông qua bằng chứng Groth16 ZK.
        \item Hàm \texttt{execute()}: thực thi lệnh chuyển tài sản thông thường sau khi xác minh ZK.
        \item Các hàm \texttt{executeStealthTransfer()}, \texttt{executeERC20StealthTransfer()}, \texttt{executeERC721StealthTransfer()}: thực thi chuyển khoản kèm phát sự kiện \texttt{Announcement}.
    \end{itemize}

    \item \textbf{\texttt{SocialRecoveryModule}:} Hợp đồng độc lập per-account, quản lý vòng đời phục hồi với ba hàm chính: \texttt{proposeRecovery()}, \texttt{approveRecovery()}, \texttt{executeRecovery()}, liên kết trực tiếp với chỉ mục lá trong Merkle Tree thông qua \texttt{mappedIndex}.

    \item \textbf{\texttt{ERC5564Announcer}:} Phát sự kiện \texttt{Announcement(schemeId, stealthAddress, caller, ephemeralPubKey, metadata)} theo chuẩn ERC-5564~\cite{erc5564}, phục vụ quá trình quét địa chỉ stealth của người nhận.

    \item \textbf{\texttt{StealthPaymaster}:} Hợp đồng Paymaster theo ERC-4337, tài trợ phí Gas cho mọi \texttt{UserOperation} nhằm giải quyết bài toán Cold Start của Stealth Account mới tạo.

    \item \textbf{\texttt{StealthSpendVerifier}} và \textbf{\texttt{SMTUpdateVerifier}:} Hai hợp đồng xác minh bằng chứng Groth16 được xuất tự động từ mạch Circom thông qua SnarkJS. \texttt{StealthSpendVerifier} dùng cho mạch chi tiêu (spend), \texttt{SMTUpdateVerifier} dùng cho mạch cập nhật cây.
\end{itemize}

\subsection{Tầng Relayer Server (Off-chain Layer)}
\label{subsec:server_layer}

Máy chủ Node.js đóng vai trò là bộ điều phối trung tâm, chịu trách nhiệm duy trì trạng thái cây và tổng hợp \texttt{UserOperation}. Tầng này hoạt động hoàn toàn ngoài chuỗi nhằm giảm chi phí tính toán Gas và tăng tốc độ phản hồi:

\begin{enumerate}
    \item \textbf{Quản lý Merkle Tree:} Server duy trì một bản sao \texttt{PoseidonSparseMerkleTree} hoàn chỉnh cục bộ, được khởi tạo từ tệp \texttt{leaves.json} và đồng bộ với trạng thái on-chain khi khởi động. Cập nhật cây sử dụng thuật toán sparse với độ phức tạp $O(d)$, chỉ tính lại các nút trên đường từ lá đến gốc.

    \item \textbf{Sinh bằng chứng ZK:} Server gọi thư viện SnarkJS để sinh bằng chứng Groth16 cho cả hai mạch: mạch chi tiêu (\texttt{stealth\_spend}) và mạch cập nhật cây (\texttt{smt\_update}). Các tệp \texttt{.wasm} và \texttt{.zkey} được nạp sẵn trong bộ nhớ để giảm trễ.

    \item \textbf{Tổng hợp UserOperation:} Sau khi có bằng chứng ZK, server xây dựng \texttt{UserOperation} hoàn chỉnh và gọi \texttt{entryPoint.handleOps()} trực tiếp thông qua ví quản lý (signer) của relayer.

    \item \textbf{Cơ chế Rollback:} Nếu giao dịch on-chain thất bại sau khi lá đã được chèn vào cây cục bộ, hàm \texttt{rollbackLastLeaf()} được gọi để hoàn tác thay đổi và duy trì tính nhất quán.
\end{enumerate}

\subsection{Tầng Client (Frontend Layer)}
\label{subsec:client_layer}

Giao diện React.js + Ethers.js v6 đảm nhiệm tất cả phép tính mật mã phía người dùng, đảm bảo khóa bí mật không bao giờ rời khỏi trình duyệt của người dùng:

\begin{itemize}
    \item Sinh và lưu trữ cặp khóa (scanPriv, spendPriv) cục bộ trong bộ nhớ phiên.
    \item Tính toán ECDH shared secret và dẫn xuất địa chỉ stealth.
    \item Quét sự kiện \texttt{Announcement} từ blockchain để phát hiện tài sản đến.
    \item Ủy quyền sinh bằng chứng ZK và gửi \texttt{UserOperation} cho server Relayer.
\end{itemize}

% ═══════════════════════════════════════════════════════════════════════════════
\section{Thiết kế Chi tiết Các Thành phần}
\label{sec:detailed_design}
% ═══════════════════════════════════════════════════════════════════════════════

\subsection{Thiết kế luồng Tạo Ví và Đăng ký Lá}
\label{subsec:create_wallet_flow}

Khi người dùng khởi tạo ví stealth mới, hệ thống thực hiện tuần tự các bước sau:

\begin{enumerate}
    \item \textbf{[Client]} Sinh cặp khóa ngẫu nhiên:
    \[
        (\texttt{scanPriv}, v) \xleftarrow{\$} \mathbb{Z}_n^*, \quad V = v \cdot G
    \]
    \[
        (\texttt{spendPriv}, x) \xleftarrow{\$} \mathbb{F}_p
    \]

    \item \textbf{[Client]} Tính giá trị lá Merkle:
    \[
        k = \text{Poseidon}(x)
    \]
    và địa chỉ danh tính:
    \[
        \texttt{addr} = \text{computeAddress}(\texttt{spendPub})
    \]

    \item \textbf{[Client $\to$ Server]} Gửi yêu cầu \texttt{POST /leaves} với payload \(\{k, \texttt{addr}\}\).

    \item \textbf{[Server]} Kiểm tra đồng bộ (\texttt{isTreeSynced}), kiểm tra trùng lá (\texttt{hasLeaf(k)}):
    \begin{itemize}
        \item Nếu lá đã tồn tại: trả về index và indexCommitment hiện có.
        \item Nếu chưa: gán chỉ mục \(i = \texttt{nextIndex}\), lấy Merkle proof của slot trống, chèn lá vào cây.
    \end{itemize}

    \item \textbf{[Server $\to$ Blockchain]} Sinh bằng chứng Groth16 SMT Update và gọi:
    \[
        \texttt{contract.updateRoot}(\texttt{newRoot}, k, i, \texttt{auth})
    \]
    Nếu giao dịch on-chain thất bại, \texttt{rollbackLastLeaf()} được gọi.

    \item \textbf{[Server $\to$ Client]} Trả về \(\{i, \texttt{newRoot}, \texttt{txHash}\}\).

    \item \textbf{[Client]} Tính:
    \[
        \texttt{IndexHash} = \text{Poseidon}(i, 0)
    \]
    và lưu bộ meta-address đầy đủ: \((\texttt{scanPriv}, \texttt{spendPriv}, V, \texttt{spendPub}, i, \texttt{IndexHash})\).
\end{enumerate}

\begin{figure}[H]
    \centering
    \includegraphics[width=1\textwidth]{SOICT_DATN_NGUYEN_VAN_AN_20224910 (1)/Hinhve/create_wallet_sq_diagram.png}
    \caption{Biểu đồ tuần tự UML -- Luồng tạo ví và đăng ký lá Merkle}
    \label{fig:sd_create_wallet}
\end{figure}

\subsection{Thiết kế luồng Gửi Tài sản Stealth}
\label{subsec:send_flow}

Luồng gửi tài sản hoàn chỉnh từ một Stealth Account đã tồn tại bao gồm 9 bước, được tổng hợp trong Bảng~\ref{table:send_flow} và minh họa trong Hình~\ref{fig:sd_send_stealth}:

\begin{table}[H]
\centering
\resizebox{\textwidth}{!}{%
\begin{tabular}{|c|l|l|p{7cm}|}
\hline
\textbf{Bước} & \textbf{Bên thực hiện} & \textbf{Hành động} & \textbf{Chi tiết} \\ \hline
1 & Client & Tra cứu Meta-address & Đọc \texttt{scanPub}, \texttt{spendPub}, \texttt{IndexHash} của người nhận (từ ENS hoặc nhập tay) \\ \hline
2 & Client & Tính ECDH Shared Secret & $S = r \cdot V_{\text{Bob}}$; $h = \text{keccak256}(S.x \| S.y)$ \\ \hline
3 & Client & Tính indexCommitment & $C = \text{Poseidon}(\text{IndexHash} + h)$ \\ \hline
4 & Client & Tính địa chỉ stealth & $A_{\text{stealth}} = \text{Factory.getAddress}(C)$ \\ \hline
5 & Client $\to$ Server & Gửi yêu cầu chi tiêu & \texttt{POST /api/spend-zk-proof} với \texttt{spendPriv}, \texttt{sharedSecretHash}, \texttt{senderIndexCommitment} \\ \hline
6 & Server & Sinh bằng chứng ZK & \texttt{generateSpendProof()} gọi SnarkJS với mạch \texttt{stealth\_spend} \\ \hline
7 & Server & Xây dựng UserOperation & \texttt{callData} = \texttt{executeStealthTransfer()}, \texttt{signature} = ABI-encode(auth, indexCommitment) \\ \hline
8 & Server & Debug xác minh on-chain & Gọi \texttt{debugVerifyProof()} trên hợp đồng trước khi gửi \\ \hline
9 & Server $\to$ Blockchain & Gửi handleOps & \texttt{entryPoint.handleOps([userOp], beneficiary)} \\ \hline
\end{tabular}
}
\caption{Luồng xử lý giao dịch stealth chi tiết}
\label{table:send_flow}
\end{table}

\begin{figure}[H]
    \centering
    \includegraphics[width=1\textwidth]{SOICT_DATN_NGUYEN_VAN_AN_20224910 (1)/Hinhve/send_stealth_payment.png}
    \caption{Biểu đồ tuần tự UML -- Luồng gửi tài sản stealth (ETH/ERC-20/ERC-721)}
    \label{fig:sd_send_stealth}
\end{figure}

\subsection{Thiết kế luồng Quét Tài sản Stealth}
\label{subsec:scan_flow}

Sau khi nhận tài sản qua giao thức Stealth Address, người nhận (Bob) cần quét blockchain để phát hiện các tài khoản thuộc sở hữu của mình. Quy trình quét được mô tả trong Hình~\ref{fig:sd_scan_stealth}:

\begin{enumerate}
    \item \textbf{[Client]} Truy vấn toàn bộ sự kiện \texttt{Announcement} từ hợp đồng \texttt{ERC5564Announcer} trong phạm vi block được chỉ định.

    \item \textbf{[Client]} Với mỗi sự kiện, tái tạo Shared Secret bằng khóa xem:
    \[
        S' = v_{\text{Bob}} \cdot R, \quad h' = \text{keccak256}(S'.x \,\|\, S'.y)
    \]

    \item \textbf{[Client]} Tính \texttt{indexCommitment} và địa chỉ stealth kỳ vọng:
    \[
        C' = \text{Poseidon}(\text{IndexHash}_{\text{Bob}} + h'), \quad A_{\text{expected}} = \text{Factory.getAddress}(C')
    \]

    \item \textbf{[Client]} So sánh \(A_{\text{expected}}\) với \texttt{stealthAddress} trong sự kiện. Nếu khớp: ghi nhận tài khoản và lưu \(h'\) phục vụ sinh ZK proof sau này.
\end{enumerate}

\begin{figure}[H]
    \centering
    \includegraphics[width=0.95\textwidth]{SOICT_DATN_NGUYEN_VAN_AN_20224910 (1)/Hinhve/scan_stealth_payment.png}
    \caption{Biểu đồ tuần tự UML -- Luồng quét và phát hiện tài sản stealth}
    \label{fig:sd_scan_stealth}
\end{figure}

\subsection{Thiết kế Cơ sở Dữ liệu Off-chain}
\label{subsec:offchain_db}

Server sử dụng tệp JSON (\texttt{leaves.json}) làm cơ sở dữ liệu ngoài chuỗi với cấu trúc như sau:

\begin{verbatim}
{
  "nextIndex": <number>,
  "leafMap": {
    "<index>": "0x<commitment_hex>",
    ...
  },
  "tree": {
    "depth": 20,
    "nodes": { ... }   // snapshot các nút non-zero
  }
}
\end{verbatim}

Mỗi trường có vai trò cụ thể:
\begin{itemize}
    \item \textbf{\texttt{nextIndex}:} Bộ đếm tăng đơn điệu, không bao giờ giảm kể cả sau khi xóa lá. Đảm bảo mỗi vị trí chỉ được gán một lần.
    \item \textbf{\texttt{leafMap}:} Ánh xạ thưa (sparse) từ chỉ mục sang giá trị cam kết lá $k$. Chỉ lưu các lá đang hoạt động (không lưu slot đã xóa).
    \item \textbf{\texttt{tree.nodes}:} Snapshot các nút nội bộ non-zero của cây Poseidon. Cho phép khởi động nhanh mà không cần tái tính toán toàn bộ cây từ các lá.
\end{itemize}

Thiết kế này giúp thời gian khởi động lại server từ $O(N \cdot d)$ (tái xây dựng đầy đủ) giảm xuống $O(1)$ (nạp snapshot).

\subsection{Thiết kế luồng Phục hồi Xã hội}
\label{subsec:recovery_flow}

Luồng phục hồi Social Recovery diễn ra qua bốn giai đoạn chính (Hình~\ref{fig:sd_social_recovery}):

\begin{enumerate}
    \item \textbf{Giai đoạn 1 -- Thiết lập Guardian (Chủ sở hữu):}  Triển khai hợp đồng \texttt{SocialRecoveryModule} với danh sách guardian và ngưỡng phê duyệt $t$. Liên kết \texttt{mappedIndex} = chỉ mục lá trong Merkle Tree tương ứng với tài khoản stealth cần bảo vệ.

    \item \textbf{Giai đoạn 2 -- Đề xuất Phục hồi (Guardian 1):}  Guardian 1 gọi \texttt{proposeRecovery(newRoot, newLeaf)}. Trước đó, server tính:
    \begin{align*}
        \texttt{oldLeaf} &= \texttt{leafMap}[i] \\
        \texttt{siblings} &= \texttt{createMerkleProof}(i) \\
        \texttt{newRoot} &= f_{\text{merkle}}(\texttt{newLeaf}, i, \texttt{siblings})
    \end{align*}
    và sinh bằng chứng ZK SMT Update qua endpoint \texttt{POST /leaves/recovery-proof}.

    \item \textbf{Giai đoạn 3 -- Bỏ phiếu Phê duyệt (Guardians 2..t):}  Các guardian còn lại gọi \texttt{approveRecovery(reqId)} để bỏ phiếu. Hợp đồng đảm bảo mỗi guardian chỉ bỏ phiếu một lần trên mỗi yêu cầu.

    \item \textbf{Giai đoạn 4 -- Thực thi Phục hồi:}  Khi \texttt{req.approvals} $\geq t$, gọi \texttt{executeRecovery(reqId, auth)}. Hợp đồng \texttt{SocialRecoveryModule} chuyển tiếp lệnh cập nhật cây với bằng chứng ZK:
    \[
        \texttt{treeManager.updateRoot(newRoot, newLeaf, mappedIndex, auth)}
    \]
    Hợp đồng \texttt{SMTUpdateVerifier} xác minh bằng chứng Groth16 trước khi cập nhật gốc on-chain.

    \item \textbf{Giai đoạn 5 -- Đồng bộ Server:}  Sau khi giao dịch on-chain được xác nhận, server đồng bộ cây cục bộ qua \texttt{POST /leaves/sync-recovery}.

    \item \textbf{Giai đoạn 6 -- Khôi phục Quyền chi tiêu:}  Chủ sở hữu mới dùng \texttt{spendPriv\_mới} để nhập ví và chi tiêu từ tài khoản stealth.
\end{enumerate}

\begin{figure}[H]
    \centering
    \includegraphics[width=1\textwidth]{SOICT_DATN_NGUYEN_VAN_AN_20224910 (1)/Hinhve/social_recovery_flow.png}
    \caption{Biểu đồ tuần tự UML -- Luồng phục hồi xã hội (Social Recovery)}
    \label{fig:sd_social_recovery}
\end{figure}

% ═══════════════════════════════════════════════════════════════════════════════
\section{Xây dựng và Triển khai Ứng dụng}
\label{sec:implementation}
% ═══════════════════════════════════════════════════════════════════════════════

\subsection{Công nghệ và Thư viện Sử dụng}

Hệ thống sử dụng các công nghệ chính được liệt kê trong Bảng~\ref{table:tech_stack}.

\begin{table}[H]
\centering
\resizebox{\textwidth}{!}{%
    \begin{tabular}{|l|l|l|l|}
        \hline
        \textbf{Thành phần} & \textbf{Công nghệ / Phiên bản} & \textbf{Vai trò} & \textbf{Phạm vi} \\ \hline
        Blockchain & Ethereum Sepolia Testnet & Môi trường triển khai và kiểm thử & On-chain \\ \hline
        Smart Contract & Solidity 0.8.28, Hardhat v3 & Phát triển, biên dịch và triển khai hợp đồng & On-chain \\ \hline
        ZK Circuits & Circom 2.0, SnarkJS 0.7 & Sinh và xác minh bằng chứng Groth16 & Off-chain + On-chain \\ \hline
        Hàm băm ZK & Poseidon (circomlibjs 0.1.7) & Băm lá và nút cây tương thích BN254 & Client + Server \\ \hline
        Frontend & React 19, Vite 7, Ethers.js v6 & Giao diện và tương tác blockchain & Client \\ \hline
        Styling & Tailwind CSS v4 & Thiết kế giao diện & Client \\ \hline
        Backend & Node.js, Express 4 & Relayer server, quản lý Merkle Tree & Off-chain \\ \hline
        Tương tác EVM & Viem v2 (Hardhat), Ethers v6 (Server) & Gọi hợp đồng và gửi giao dịch & Server + Client \\ \hline
        ENS & ETHRegistrarController Sepolia & Đăng ký tên miền .eth và resolver records & On-chain \\ \hline
    \end{tabular}
}
\caption{Bảng công nghệ và thư viện sử dụng trong hệ thống}
\label{table:tech_stack}
\end{table}

\subsection{Cấu trúc Module Hệ thống}
\label{subsec:module_structure}

\subsubsection{Module Smart Contract (\texttt{blockchain/})}

\begin{itemize}
    \item \texttt{contracts/zk/}: Hợp đồng \texttt{StealthSpendVerifier.sol} và \texttt{SMTUpdateVerifier.sol} được tự động sinh từ SnarkJS.
    \item \texttt{contracts/stealth-abstract-account/}: Bốn hợp đồng cốt lõi (\texttt{StealthAccount}, \texttt{StealthAccountFactory}, \texttt{StealthPaymaster}, \texttt{SocialRecoveryModule}).
    \item \texttt{scripts/deploy.ts}: Script triển khai tuần tự theo thứ tự phụ thuộc.
\end{itemize}

\subsubsection{Module ZK (\texttt{zk/circuits/})}

\begin{itemize}
    \item \texttt{stealth\_spend/stealth.circom}: Mạch ZK xác thực quyền chi tiêu (độ sâu 20, $\approx$ 11.7k ràng buộc R1CS).
    \item \texttt{smt/smt\_update.circom}: Mạch ZK xác thực cập nhật cây (độ sâu 20, $\approx$ 20.7k ràng buộc R1CS).
    \item Mỗi mạch đi kèm: \texttt{.zkey} (khóa chứng minh), \texttt{verification\_key.json}, và \texttt{\_js/} (wasm witness generator).
\end{itemize}

\subsubsection{Module Relayer Server (\texttt{server/})}

\begin{description}
    \item[\texttt{server.js}] Điểm khởi động, gắn kết các router và khởi tạo cây Merkle.
    \item[\texttt{services/merkleService.js}] Toàn bộ logic quản lý cây: khởi tạo, đọc/ghi lá, sinh Merkle proof, đồng bộ on-chain khi khởi động.
    \item[\texttt{services/zkService.js}] Giao diện SnarkJS: \texttt{generateSpendProof()} và \texttt{generateUpdateProof()}.
    \item[\texttt{services/contractService.js}] Gọi \texttt{contract.updateRoot()} sau khi có bằng chứng ZK.
    \item[\texttt{routes/leaves.js}] API quản lý lá: đăng ký, tìm kiếm, xóa, phục hồi.
    \item[\texttt{routes/spend.js}] API chi tiêu end-to-end: sinh ZK proof, xây dựng và gửi \texttt{UserOperation}.
    \item[\texttt{routes/tree.js}] API đọc: root hiện tại, Merkle proof theo chỉ mục.
\end{description}

\subsubsection{Module Frontend (\texttt{stealth-wallet/src/})}

\begin{description}
    \item[\texttt{stealth/crypto.js}] Tạo cặp khóa meta-address, tính ECDH, dẫn xuất địa chỉ stealth.
    \item[\texttt{stealth/zkIntegration.js}] Giao tiếp với server: đăng ký lá, tra cứu lá, tính \texttt{indexCommitment}.
    \item[\texttt{scanner/scannerEngine.js}] Quét sự kiện \texttt{Announcement} từ blockchain và lọc giao dịch thuộc sở hữu người dùng.
    \item[\texttt{services/stealthService.js}] Điều phối luồng chi tiêu: gọi \texttt{/api/spend-zk-proof}.
    \item[\texttt{services/ensRegistrationService.js}] Giao tiếp với ENS: commit, đăng ký và ghi resolver records.
    \item[\texttt{hooks/useStealthWallet.js}] React hook tổng hợp: tạo ví, nhập ví, quét, gửi.
    \item[\texttt{components/}] Năm màn hình chính: \texttt{Wallet}, \texttt{Send}, \texttt{Assets}, \texttt{EnsRegister}, \texttt{Recovery}.
\end{description}

\subsection{Giao diện Người dùng}
\label{subsec:ui}

Giao diện được thiết kế theo phong cách hiện đại, sáng màu (light mode) với thiết kế dạng thẻ (card-based), sử dụng Tailwind CSS v4. Các thành phần chính bao gồm:

\begin{itemize}
    \item \textbf{Màn hình Dashboard (Wallet):} Hiển thị meta-address (scanPub, spendPub, indexHash) và tổng số dư tài sản. Hỗ trợ quét và liệt kê các tài khoản stealth phát hiện được kèm số dư (ETH, ERC-20, ERC-721). Tích hợp giao diện gửi tài sản inline thông qua bằng chứng ZK ngay trong mỗi thẻ tài khoản.

    \item \textbf{Màn hình Transfer (Send):} Giao diện chuyển tài sản trực tiếp (Direct Transfer) từ ví chính đến địa chỉ stealth của người nhận khác, hỗ trợ tự động điền khóa công khai thông qua tra cứu ENS.

    \item \textbf{Màn hình Security (Recovery):} Cấu hình danh sách guardian và ngưỡng phê duyệt. Màn hình guardian riêng để đề xuất, bỏ phiếu và thực thi phục hồi tài khoản stealth.

    \item \textbf{Màn hình ENS Domain:} Giao diện hai bước Commit $\to$ Register với bộ đếm thời gian chờ 60 giây bắt buộc. Tự động ghi ba bản ghi stealth sau khi đăng ký tên miền thành công.
\end{itemize}

Mỗi bước có thao tác dài (sinh bằng chứng ZK, ký giao dịch) đều có hiển thị trạng thái tiến trình chi tiết (ví dụ: ``Generating ZK proof...'', ``Submitting UserOperation...'', ``Waiting for confirmation...'') để người dùng theo dõi quá trình thực thi.

% ═══════════════════════════════════════════════════════════════════════════════
\section{Kiểm thử và Đánh giá Hệ thống}
\label{sec:testing}
% ═══════════════════════════════════════════════════════════════════════════════

\subsection{Kiểm thử Chức năng}
\label{subsec:functional_testing}

Hệ thống được kiểm thử theo các kịch bản chức năng sau trên mạng Sepolia Testnet:

\begin{enumerate}
    \item \textbf{Kiểm thử Tạo ví và Đăng ký Lá:}
    \begin{itemize}
        \item Tạo ví mới: xác nhận lá $k = \text{Poseidon}(x)$ được chèn đúng vào cây và gốc Merkle on-chain khớp với gốc server.
        \item Nhập ví đã tồn tại: xác nhận hệ thống phát hiện lá trùng và trả về chỉ mục cũ thay vì chèn mới.
        \item Rollback khi on-chain thất bại: xác nhận \texttt{leaves.json} hoàn nguyên về trạng thái trước khi chèn.
    \end{itemize}

    \item \textbf{Kiểm thử Quét Địa chỉ Stealth:}
    \begin{itemize}
        \item Quét đúng: xác nhận chỉ các tài khoản thuộc về người dùng (có \texttt{scanPriv} tương ứng) được nhận diện.
        \item Bỏ qua đúng: xác nhận các giao dịch không thuộc về người dùng bị lọc ra.
    \end{itemize}

    \item \textbf{Kiểm thử Giao dịch Chi tiêu:}
    \begin{itemize}
        \item Giao dịch hợp lệ: bằng chứng ZK đúng với \texttt{root}, \texttt{indexCommitment} khớp -- giao dịch được thực hiện thành công.
        \item Giao dịch với bằng chứng sai: \texttt{debugVerifyProof()} trả về \texttt{false}, giao dịch bị từ chối trước khi gửi.
        \item Giao dịch không đủ số dư: server kiểm tra số dư ETH và trả về lỗi \texttt{INSUFFICIENT\_BALANCE} trước khi sinh bằng chứng.
        \item Giao dịch ERC-20: xác nhận \texttt{executeERC20StealthTransfer()} gọi đúng ABI và phát sự kiện Announcement kèm metadata.
    \end{itemize}

    \item \textbf{Kiểm thử Phục hồi Xã hội:}
    \begin{itemize}
        \item Luồng đề xuất - bỏ phiếu - thực thi khi đủ ngưỡng: xác nhận gốc Merkle on-chain cập nhật đúng giá trị mới.
        \item Bỏ phiếu trùng: hợp đồng từ chối lần bỏ phiếu thứ hai của cùng một guardian.
        \item Thực thi khi chưa đủ ngưỡng: hợp đồng revert với thông báo \texttt{not enough approvals}.
    \end{itemize}

    \item \textbf{Kiểm thử ENS:}
    \begin{itemize}
        \item Đăng ký trước thời gian chờ 60 giây: hợp đồng từ chối.
        \item Đăng ký đúng quy trình: tên miền được tạo và ba bản ghi stealth được ghi thành công vào resolver.
    \end{itemize}
\end{enumerate}

\subsection{Đánh giá Hiệu năng}
\label{subsec:performance}

\subsubsection{Thời gian Sinh Bằng chứng ZK}

Quá trình sinh bằng chứng Groth16 được thực hiện phía server. Bảng~\ref{table:proof_time} tổng hợp thời gian đo được trên môi trường thử nghiệm (CPU: Intel Core i7, RAM: 16 GB):

\begin{table}[H]
\centering
\resizebox{\textwidth}{!}{%
\begin{tabular}{|l|c|c|c|}
\hline
\textbf{Mạch} & \textbf{Số ràng buộc R1CS} & \textbf{Thời gian trung bình} & \textbf{Thời gian tối đa} \\ \hline
\texttt{stealth\_spend} (depth=20) & 11,749 & 0.2--0.3 giây & 0.5 giây \\ \hline
\texttt{smt\_update} (depth=20) & 20,781 & 0.2--0.4 giây & 0.5 giây \\ \hline
\end{tabular}%
}
\caption{Thời gian sinh bằng chứng Groth16 trên server}
\label{table:proof_time}
\end{table}

Được thực hiện bất đồng bộ phía server, thời gian này không ảnh hưởng đến trải nghiệm giao diện người dùng.

\subsubsection{Chi phí Gas trên Blockchain}

Bảng~\ref{table:gas_cost} trình bày ước tính chi phí Gas cho từng thao tác trên Sepolia Testnet:

\begin{table}[H]
\centering
\resizebox{\textwidth}{!}{%
\begin{tabular}{|l|c|l|}
\hline
\textbf{Thao tác} & \textbf{Gas tiêu thụ (ước tính)} & \textbf{Ghi chú} \\ \hline
Xác minh Groth16 (spend) & $\approx$245,000 & Không đổi theo kích thước cây \\ \hline
Xác minh SMT Update & $\approx$270,000 & Không đổi theo kích thước cây \\ \hline
Triển khai \texttt{StealthAccount} & $\approx$450,000 & Chỉ thực hiện một lần \\ \hline
Chuyển ETH có Announcement & $\approx$80,000 & Thêm sự kiện ERC-5564 \\ \hline
Chuyển ERC-20 có Announcement & $\approx$95,000 & Tùy thuộc vào token contract \\ \hline
Đăng ký ENS (1 năm) & $\approx$300,000 & Không tính phí ETH đăng ký \\ \hline
\end{tabular}%
}
\caption{Ước tính chi phí Gas trên Ethereum Sepolia}
\label{table:gas_cost}
\end{table}

\textbf{Nhận xét quan trọng:} Chi phí xác minh Groth16 duy trì ở mức cố định $O(1)$ không phụ thuộc vào kích thước cây Merkle ($2^{20}$ lá). Đây là ưu điểm cốt lõi của hệ thống chứng minh ZK-SNARK so với cách tiếp cận xác minh trực tiếp Merkle Tree trên EVM, vốn yêu cầu chi phí $O(d)$ theo độ sâu cây.

\subsubsection{Tính Đồng bộ và Độ tin cậy}

\begin{itemize}
    \item \textbf{Cơ chế kiểm tra tại khởi động:} Mỗi lần server khởi động, hệ thống so sánh gốc cây cục bộ với gốc on-chain. Nếu phát hiện lệch (\texttt{ROOT MISMATCH}), mọi thao tác chèn lá mới bị từ chối cho đến khi đồng bộ lại.

    \item \textbf{Cơ chế Rollback lạc quan:} Nếu giao dịch on-chain thất bại sau khi cây cục bộ đã được cập nhật, hệ thống tự động gọi \texttt{rollbackLastLeaf()} để khôi phục về trạng thái trước, ngăn chặn mất đồng bộ vĩnh viễn.

    \item \textbf{Idempotency:} Gọi \texttt{deployFor(indexCommitment)} nhiều lần trên Factory luôn trả về cùng địa chỉ hợp đồng đã tồn tại thay vì tạo mới, đảm bảo an toàn khi có lỗi mạng và retry.
\end{itemize}

\end{document}


